\clearpage
\fancyhead[L]{\small PRÄCESSION.}
\SourceLeaf{566}
\begin{center}\small Göttingische gelehrte Anzeigen. Stück 19. Seite 188..189. 1814 Januar 31.\\[3pt]\rule{.72\linewidth}{.3pt}\end{center}

Die königl. Academie der Wissenschaften in Berlin hatte die genaueste Bestimmung der Grösse der Präcession zum Gegenstande einer Preisfrage gemacht: sie hat der Arbeit des Hrn. \textsc{Bessel} darüber den Preis zuerkannt. Niemand konnte auch in der That diese Bestimmung mit glücklicherm Erfolg übernehmen, als Hr. \textsc{Bessel}, welcher seit sechs Jahren an einer vielseitigen Discussion der \textsc{Bradley}-schen Beobachtungen gearbeitet hat. Er hat zu diesem Behuf 4585 Sternbestimmungen angewandt, und wir glauben den Astronomen einen wichtigen Dienst zu leisten, wenn wir das von ihm gefundene Endresultat für eine Grösse, welche sie täglich nöthig haben, hier mittheilen.

\begin{center}\fontsize{8}{11}\selectfont\renewcommand{\arraystretch}{1.18}\begin{tabular}{@{}l@{\hspace{4pt}}l@{}}
Lunisolarpräcession $.\;.\;.\;.\;.\;.\;.\;.\;.\;.\;.\;.$ & $= 50''35330 - 0''0002435890\,(t-1800)$ \\
Beobachtete allgemeine Präcession $.\;.\;.\;.\;.\;.$ & $50.18924 + \mathrm{c}.0002442966\,(t-1800)$ \\
Constante bei der Präcession in Gerader Aufsteigung & $46.01058 + 0.0003590677\,(t-1800)$ \\
Constante bei der Präcession in Declination $.\;.$ & $20.04966 - 0.0002135621\,(t-1800)$ \\
\end{tabular}\end{center}

Von der von \textsc{Triesnecker} bemerkten Erscheinung, dass die Bewegungen an einigen Punkten des Himmels auf eine ungleiche Präcession deuten sollen, hat Herr \textsc{Bessel} keine deutliche Spur bemerkt.
